\subsection{Repair diagnostics: fixed protocol and equal order weights}
The versioned protocol \texttt{group-subset-v1} now drives training and EDA.
Historical artifacts are preserved separately. The already observed holdout is
reused only for reproduction and descriptive evaluation, not new tuning.
The subset contains 5999 of 180519 rows and 2113 of 65752 orders.
The largest absolute category-proportion difference across the audited features
and target is 0.0227 (Customer Segment).
These marginal checks do not establish unbiased sampling.

For the selected tree B, equal-order evaluation over 423 orders gives
Accuracy=0.7258, Precision=0.8554, Recall=0.6068,
F1=0.7100; TN=165, FP=24, FN=92, TP=142.
The fixed rule averages item probabilities per order and predicts positive at
0.5. It differs from row weighting and is not a measured operational benefit.

The controlled calendar ablation uses the same training folds and A-selected
hyperparameters for both flows. Scores below are training-CV validation scores;
the CSV's test-prefixed names are sklearn fold-validation names, not holdout scores.
\begin{center}
\begin{tabular}{lrr}
Algorithm & A CV F1 & B CV F1 \\
knn & 0.6668 & 0.6714 \\
tree & 0.6741 & 0.6904 \\
\end{tabular}
\end{center}
The tree's independently tuned B score also reflects a different depth choice;
the fixed-parameter comparison separates that change from adding calendar inputs.
Because parameters were selected on these folds, these are descriptive diagnostics
with selection optimism, not causal effects or independent generalization evidence.

Standard Class has 248 FN among 254 positive rows (FNR 0.9764), whereas
First/Second Class have 10/40 FP among 10/40 negative rows respectively.
The selected tree trades high precision for many missed Standard Class delays.
Use the exported rules and denominators to describe its behavior; do not infer
causality or adjust this model against the known holdout.
Canceled orders retain source labels; target zero does not exclusively denote
successful on-time delivery. Operational label redesign requires a new protocol.
